%% =====================================================================
%%  T-25-02  The Free Lunch
%%  -------------------------------------------------------------------
%%  One idea per file. Core is mandatory. Slots in canonical order.
%%  Max 700 words. No layout commands. No filler. New source material
%%  for this entry goes in sources/<BOOK>/T-25-02.tex -- never by
%%  rewriting this file (docs/RULES.md R0.1).
%%  =====================================================================
%% @kind: topic
%% @id: T-25-02
%% @title: The Free Lunch
%% @subtitle: What obligation does to a price
%% @chapter: c25
%% @order: 20
%% @status: stable
%% @sources: B3
%% @updated: 2026-09-19
%% @allow-rewrite: B3 ingest 2026-09-19 -- committed scaffold replaced by the written entry (planned -> stable lifecycle)

\topic{T-25-02}{The Free Lunch}{What obligation does to a price}

\begin{core}
There are no bargains: what arrives without a price arrives with a hook. The free lunch creates obligation (the reciprocity lever of T-19-01), lowers your guard (suspicion is what the gift is for), and marks you as the kind of person who can be bought for the price of a meal. Despise the free lunch --- pay full price, and keep your judgement unencumbered.
\end{core}
\begin{mechanism}
Unpaid-for benefits are the cheapest opening move in the entire craft because they carry three payloads at once. First, obligation: the gift opens an account whose balance the giver sets later (T-19-01 is the mechanism; this entry is its market price). Second, disarmament: the tested way to enter a fortress is as a friend bearing gifts --- suspicion is precisely what the gift exists to suspend, and the flag of trade is run up so the flag of your rising guard can come down. Third, classification: accepting free things teaches the giver's network that you transact in favours, which converts you from a person who pays into a person who can be paid --- the cheapest kind of asset there is. The Greeks' warning at Troy was already old: fear Greeks bearing gifts. The trap's genius is asymmetry --- the gift costs the giver little and obligates you at the moment of your choosing, not theirs.
\end{mechanism}
\begin{conditions}
Strongest where relationships compound --- ongoing business, courts, communities --- because the account opened by a meal accrues. Also strongest against the proud, who believe themselves above such leverage precisely because they accept generously. It weakens in one-shot markets with settled prices, and against practitioners who reciprocate instantly, which zeroes the account.
\end{conditions}
\begin{application}
  \item Pay asking prices without ceremony for things that matter. The premium you pay is the audit you skip.
  \item When a gift arrives, price its hook before accepting: what does the giver gain by holding an account on you?
  \item If you accept --- and business makes you accept --- reciprocate fast, in kind, visibly. Close the account the same week.
  \item In your own dealings, charge what you are worth. Discounting is self-selling on instalments.
  \item Track who buys, who gives free, and who discounts around you. Each is running a different strategy on your ledger.
\end{application}
\begin{feedback}
  \item Working: the giver's requests arrive later, larger, and prefaced with \emph{after all I did}. The account is being called.
  \item Working: your paid-full-price purchases stay clean --- no ask follows them, ever.
  \item Failing: you notice you are defending the giver to others. The gift is now spending you.
  \item Failing: everything you buy discounted comes with advice attached. The discount was the product.
\end{feedback}
\begin{failure}
The rule calculates into paranoia --- a person who reads every kindness as a hook becomes unreachable by genuine generosity and ends up poorer in the only currency that compounds. It also has a real cost in reciprocity: refusing all gifts refuses the social glue, and some gifts are just gifts. The skill is in pricing speed --- the hook is usually visible in what the giver does in the week after you accept.
\end{failure}
\begin{countermeasures}
  \item In your own buying, be the counter-example: pay real prices, on time, without negotiation theatre. It keeps your own judgement expensive to rent.
  \item When someone insists on giving, ask yourself what they are buying --- and notice the answer is usually access, gratitude, or a precedent.
  \item Teach your team the flag trick: the gift before the pitch is the pitch. Reward the members who bill properly.
  \item Keep one standing rule visible: you reciprocate within the week. Known rules deter the account-openers before the meal is ordered (T-04-02).
\end{countermeasures}

\sources{B3}
\seesources
\seealso{T-19-01, T-04-02, T-25-01}
